\section{Ordinary projectors, finite freeness, and coefficient base change}
\label{sec:ordinary-structure-v42}

\status{This chapter reconstructs the original generic endpoints
\texttt{OrdFree.ordStruct}, \texttt{OrdFree.ordFree}, and
\texttt{OrdBCMapOrd.ordBC}, including their elementary input producers.
The hypotheses are precisely the original finite-level, stabilizer, and
continuous-character hypotheses on a supplied Frame.  The coefficient
rings are arbitrary complete Noetherian local rings with finite residue
fields; reducedness, flatness of a coefficient map, and characteristic
zero are not required.  The construction uses the actual function spaces,
double-coset operators, and coefficientwise maps.  Applying it to the
original arithmetic data still uses the separately constructed Frame and
its finite-level and weight inputs.  Classical control, arithmetic fibers,
and the same-ideal Hida endpoints are not conclusions of this chapter.}

\subsection{The actual spaces and their hypotheses}

Write $G=G_a\times G_p$ and let $\operatorname{glob}:\Gamma\to G$
and $U_b\leq G_a$ be the supplied Frame.  Its Iwahori datum consists of
subgroups $N\leq N_f$, $T\leq T_f$, decreasing subgroups
$T_c(c)\leq T$ and $N^-(c)$, a submonoid $P\subset T_f$, and
$z\in P$.  We retain every original Iwahori axiom.  The ones used
below have the following mathematical content:

\begin{enumerate}
\item $T_f$ is commutative and normalizes $N_f$; $T$ normalizes
both $N$ and every $N^-(c)$.
\item For $a\in P$, $aNa^{-1}\subset N$ has finite index, and
$a^{-1}N^-(c)a\subset N^-(c)$.  In addition,
$z^{-1}N^-(c)z\subset N^-(c+1)$.
\item For $y\in N^-(c)$ and $x\in N$, there are
$x'\in N$, $t\in T_c(c)$, and $y'\in N^-(c)$ with
$yx=x'ty'$.
\item If $xty=1$, with $x\in N_f$, $t\in T_f$, and
$y\in N^-(0)$, then $x=t=1$ (and hence $y=1$).
\item The indices $[N^-(0):N^-(c)]$ and $[T:T_c(c)]$ are
finite.  For every $a\in P$, some $h\geq0$ satisfies
$z^h a^{-1}\in P$.
\end{enumerate}

Put $I(c)=\langle N,T,N^-(c)\rangle$ and
$I_1(c)=\langle N,T_c(c),N^-(c)\rangle$.  Let $(A,\mathfrak m)$
be a complete Noetherian local ring with finite residue field, let
$d\geq0$, and fix
\[
 \tau:U_b\longrightarrow\operatorname{GL}_d(A),\qquad
 \theta:T\longrightarrow A^\times,\qquad U\leq U_b.
\]
The original continuous-character condition is
\begin{equation}\label{eq:ost-contchar}
 \forall k\geq0\ \exists c\geq0\quad
             \theta(t)-1\in\mathfrak m^k\quad(t\in T_c(c)).
\end{equation}
The finite-level condition is finiteness of
$\operatorname{glob}(\Gamma)\backslash G/(U\times I(0))$,
with the factorization order
$g=\operatorname{glob}(\gamma)x(u,b)$.  The stabilizer condition
is the original \texttt{Free}: whenever
\begin{equation}\label{eq:ost-stabilizer}
 \operatorname{glob}(\gamma)g=g(u,b),\quad u\in U,\ b\in I(0),
\end{equation}
one has $b\in Z(G_p)\cap T$ and $\theta(b)\tau(u)=1$.
This is a condition on stabilizers; it does not presume module freeness.

Let $\mathcal F_A$ be the module of all functions $G\to A^d$.
For $a\in G_p$ write $r_a f(g)=f(g(1,a))$.  The actual space
$M_A=\mathrm{Sp}_A(\tau,\theta,U)$ consists of functions satisfying
\begin{align}
 f(\operatorname{glob}(\gamma)g)&=f(g), &
 r_n f&=f\quad(n\in N),\\
 r_t f&=\theta(t)^{-1}f\quad(t\in T), &
 \tau(u)f(g(u,1))&=f(g)\quad(u\in U),\\
 \forall k\ \exists c\quad
 r_y f-f&\in\mathrm{valIn}_A(k)\quad(y\in N^-(c)).
 \label{eq:ost-space}
\end{align}
Here $\mathrm{valIn}_A(k)$ means that every coordinate of every
function value lies in $\mathfrak m^k$.  Write
$V_A(k)=M_A\cap\mathrm{valIn}_A(k)$.  In particular these
are value filtrations, before any identification with intrinsic
module filtrations.

For $a\in P$ choose representatives $R_a$ of the left cosets
$N/(aNa^{-1})$.  On $N$-fixed functions the original operator is
\begin{equation}\label{eq:ost-ua}
                 U_a f=\sum_{n\in R_a}r_{na}f.
\end{equation}
It agrees with the source's sum using quotient representatives.
The ordinary space is the literal source submodule
\begin{equation}\label{eq:ost-ord-definition}
 O_A=M_A\cap\bigcap_{k\geq0}
       \sum_{q\geq1}\ker\bigl(U_z^q-1:
                         \mathcal F_A\to\mathcal F_A/
                                      \mathrm{valIn}_A(k)\bigr).
\end{equation}
The sum notation in this definition is made concrete in OST10.

\subsection{Finite coordinates without changing the function module}

% BEGIN OST1
\paragraph{OST1. Iwahori factorization and finite representatives at every depth.}
\label{claim:OST1}
One has
\[
 I(c)=NTN^-(c)=N^-(c)TN,\qquad
 I_1(c)=NT_c(c)N^-(c),\qquad I(c)\cap N_f=N.
\]
Moreover $[I(0):I(c)]$ and $[I(0):I_1(c)]$ are finite.
Consequently the finite-level hypothesis gives, for every $c$, a
finite list $x_1,\ldots,x_s$ of distinct double-coset representatives
such that every $g\in G$ can be written
\begin{equation}\label{eq:ost-cosets}
             g=\operatorname{glob}(\gamma)x_i(u,b),
                       \quad u\in U,\quad b\in I(c).
\end{equation}
The index $i$ is unique.
\begin{proof}
More generally let $T_c(c)\leq S\leq T$.  The set $NSN^-(c)$
is closed under multiplication: in a product $xty\,x't'y'$, use
$yx'=x''s y''$ with $s\in T_c(c)$, move torus elements through
$N$ and $N^-(c)$ by their normalization, and collect the three
factors.  It contains the three subgroups, and it is closed under
inversion, since $(xty)^{-1}=y^{-1}t^{-1}x^{-1}$ is a product
of three elements of that same multiplicatively closed set.
It is therefore their generated subgroup.  This proves the forward
factorizations; inversion proves the reversed one.  If $g=xty$
also belongs to $N_f$, apply the uniqueness axiom to
$(g^{-1}x)ty=1$.  It gives $g=x\in N$.

The reversed factorization gives $I(0)=N^-(0)I(c)$, and
$N^-(c)\subset I(c)$.  Thus the finite coset set
$N^-(0)/N^-(c)$ surjects onto $I(0)/I(c)$.
Similarly $I(c)=T I_1(c)$ and $T_c(c)\subset I_1(c)$, so
$[I(c):I_1(c)]\leq[T:T_c(c)]$.
Refine a finite list at level $I(0)$ by representatives of
$I(0)/I(c)$ and discard duplicate double cosets.  This gives
\eqref{eq:ost-cosets} with the stated uniqueness.  No normality
of $I(c)$ or averaging denominator is used.
\end{proof}
% END OST1

% BEGIN OST2
\paragraph{OST2. The torus coordinate and its congruence variation.}
\label{claim:OST2}
There is a well-defined function $t_0:I(0)\to T$ characterized by
$t_0(y t n)=t$ for $y\in N^-(0)$, $t\in T$, $n\in N$.
It satisfies
\begin{gather}
 t_0(1)=1,\quad t_0(sb)=s t_0(b),\quad
 t_0(bs)=t_0(b)s,\quad t_0(bn)=t_0(b),\\
 y\in N^-(c)\quad\Longrightarrow\quad
 t_0(by)=t_0(b)s_c\quad\hbox{for some }s_c\in T_c(c).
 \label{eq:ost-torus}
\end{gather}
Here $s\in T$, $n\in N$, and $b\in I(0)$.
\begin{proof}
Existence of the reversed factorization is OST1.  If
$ytn=y't'n'$, rearranging gives an identity with successive factors
in $N_f,T_f,N^-(0)$.  Explicitly, from
$n=t^{-1}y^{-1}y't'n'$ one obtains
\[
 (nn'^{-1})(t'^{-1}t)
           \bigl(t^{-1}y'^{-1}yt\bigr)=1.
\]
The last factor belongs to $N^-(0)$ by torus normalization.
The uniqueness axiom gives $t'^{-1}t=1$.
The first four formulas now follow by moving a torus factor through
$N^-(0)$ or $N$, or multiplying two $N$ factors.

For the last formula write $b=y_0tn$.  Apply the swap axiom at
depth $c$ to $y^{-1}n^{-1}=n'sy'$ with $s\in T_c(c)$.
Then $ny=(y')^{-1}s^{-1}(n')^{-1}$.  Move $(y')^{-1}$
past $t$ and combine it with $y_0$; the resulting reversed
factorization of $by$ has torus coordinate $ts^{-1}$.
Since $T_c(c)$ is a subgroup, this is the required assertion.
\end{proof}
% END OST2

% BEGIN OST3
\paragraph{OST3. Integral sections for the actual coefficient functions.}
\label{claim:OST3}
Fix the representatives of OST1 at depth $c$.  There is an
$A$-linear map
\[
 s_c:(A^d)^s\longrightarrow M_A
\]
such that $(s_c v)(x_i)=v_i$.  If
$\theta(T_c(c))\subset1+\mathfrak m^k$, then every $s_cv$ is
$N^-(c)$-invariant modulo $\mathrm{valIn}_A(k)$.
\begin{proof}
For a factorization \eqref{eq:ost-cosets} define
\begin{equation}\label{eq:ost-section}
 (s_cv)(g)=\theta(t_0(b))^{-1}\tau(u^{-1})v_i.
\end{equation}
To verify independence of the factorization, compare
$\operatorname{glob}(\gamma)x_i(u,b)$ and
$\operatorname{glob}(\gamma_1)x_i(u_1,b_1)$.  The index is the
same by OST1.  The resulting stabilizer relation is
\[
 \operatorname{glob}(\gamma_1^{-1}\gamma)x_i
       =x_i(u_1u^{-1},b_1b^{-1}).
\]
By \eqref{eq:ost-stabilizer}, $s=b_1b^{-1}\in T$ and
$\theta(s)\tau(u_1u^{-1})=1$.  Hence
\[
 t_0(b_1)=s t_0(b),\qquad
 \tau(u_1^{-1})=\theta(s)\tau(u^{-1}).
\]
These factors cancel in \eqref{eq:ost-section}.  The formula is
linear in $v$ and gives $v_i$ at $x_i$.

Left $\Gamma$-invariance follows by changing $\gamma$.
Right $N$-invariance and the $T$-character law follow from OST2.
For $u'\in U$, replace $u$ by $uu'$ and use
$\tau(u')\tau((uu')^{-1})=\tau(u^{-1})$; this is the required
away-factor equivariance.  Finally fix $k$ and choose
$c'\geq c$ with $\theta(T_c(c'))\subset1+\mathfrak m^k$.
For $y\in N^-(c')$, the last formula of OST2 changes
\eqref{eq:ost-section} by multiplication by
$\theta(s_{c'})^{-1}$, which is congruent to one modulo
$\mathfrak m^k$.  Indeed $a^{-1}-1=-a^{-1}(a-1)$ for a unit
$a$.  All entries of $\tau(u^{-1})$ lie in $A$, so the
congruence holds for every $g$ and every coordinate.  This proves
the last condition of \eqref{eq:ost-space}.  If the character is
already trivial at depth $c$ modulo $\mathfrak m^k$, take $c'=c$.
\end{proof}
% END OST3

% BEGIN OST4
\paragraph{OST4. Finite values detect congruences at a fixed depth.}
\label{claim:OST4}
Suppose $f\in M_A$, $r_yf-f\in\mathrm{valIn}_A(k)$ for
$y\in N^-(c)$, and $f(x_i)\in(\mathfrak m^k)^d$ at all
representatives of OST1.  Then $f\in V_A(k)$.
\begin{proof}
Write $g=\operatorname{glob}(\gamma)x_i(u,b)$ and factor
$b=nty$ by OST1.  Modulo $\mathfrak m^k$, remove $y$ using
the assumed invariance, remove $t$ using its unit character,
and remove $n$ using exact $N$-invariance.  Away-factor
equivariance then expresses $f(g)$ modulo $\mathfrak m^k$ as
$\theta(t)^{-1}\tau(u^{-1})f(x_i)$.  Its entries belong to
$\mathfrak m^k$.  This proof uses precisely the integral matrix
$\tau(u^{-1})$, not an unspecified lattice in an isomorphic space.
\end{proof}
% END OST4

\subsection{The original operators and level lowering}

% BEGIN OST5
\paragraph{OST5. Double-coset algebra with no division.}
\label{claim:OST5}
On $N$-fixed vectors in any representation, the sum
\eqref{eq:ost-ua} is independent of its representatives, has
$N$-fixed image, and satisfies
\[
 U_aU_b=U_{ab},\qquad U_1=1,\qquad r_tU_a=U_a r_t
                    \quad(a,b\in P,\ t\in T).
\]
For function representations it commutes with all coefficientwise
ring maps, preserves every value filtration, and commutes with
the away action $f(g)\mapsto\tau(u)f(g(u,1))$.
\begin{proof}
Replacing $n\in R_a$ by $n a n_0a^{-1}$ changes
$r_{na}f$ into $r_{na}r_{n_0}f$, which is the same vector.
Left multiplication of $R_a$ by an element of $N$ permutes its
cosets, proving $N$-invariance.  If $R_a$ and $R_b$ are chosen,
then
\[
             \{n(a n'a^{-1}):n\in R_a,\ n'\in R_b\}
\]
is a representative set for $N/(abN(ab)^{-1})$: first reduce
modulo $aNa^{-1}$ and then conjugate the remaining factor by
$a^{-1}$ and reduce modulo $bNb^{-1}$.  Uniqueness follows
by reversing those two steps.  Expanding the two finite sums
proves the product identity.  For $a=1$ there is one coset.
Conjugation by $t\in T$ permutes the quotient because $t$
normalizes $N$ and commutes with $a$; this proves the torus
identity.  The final assertions follow term by term from the
finite sum.  No index is inverted, so these identities hold in
every characteristic and with nilpotent coefficients.
\end{proof}
% END OST5

% BEGIN OST6
\paragraph{OST6. Level preservation and strict lowering by $U_z$.}
\label{claim:OST6}
If a vector $v$ is fixed by $N$, by $T_c(c)$, and by every
$a^{-1}N^-(c)a$, then $U_av$ is fixed by $N^-(c)$.
Consequently $U_a$ preserves $I_1(c)$-invariants, and if
$d'\geq c$, then $U_z$ sends vectors fixed by
$N,T_c(c),N^-(d'+1)$ to vectors fixed by
$N,T_c(c),N^-(d')$.
\begin{proof}
Fix $y\in N^-(c)$.  For each representative $n\in R_a$,
the swap axiom gives $yn=x t y'$ with $x\in N$,
$t\in T_c(c)$, $y'\in N^-(c)$.  Choose $n'\in R_a$
with $x=n'a n_0a^{-1}$, $n_0\in N$.  Since $t$ commutes
with $a$,
\[
 yn a=n'a k_n,\qquad k_n=n_0t(a^{-1}y'a)\in I(c),
                       \qquad r_{k_n}v=v.
\]
The resulting map $n\mapsto n'$ is injective.  If two images
coincide, then
$a^{-1}n^{-1}\widetilde n a=k_n^{-1}k_{\widetilde n}$
belongs to $I(c)\cap N_f=N$ by OST1.  Thus $n$ and
$\widetilde n$ represent the same coset and are equal.
It is a permutation of the finite set, so summing gives
$r_yU_av=U_av$.

The remaining $N$ and torus invariances are OST5.
For level preservation use
$a^{-1}N^-(c)a\subset N^-(c)$.
For strict lowering apply the argument at depth $d'$,
using $T_c(d')\subset T_c(c)$ and
$z^{-1}N^-(d')z\subset N^-(d'+1)$.
\end{proof}
% END OST6

% BEGIN OST7
\paragraph{OST7. Stability of the actual automorphic space and eventual lowering.}
\label{claim:OST7}
Every $U_a$, $a\in P$, preserves $M_A$ and its value
filtration.  If $\theta(T_c(c))\subset1+\mathfrak m^k$,
then for every $f\in M_A$ there is $n_0$ such that
\begin{equation}\label{eq:ost-eventual-lowering}
 n\geq n_0,\ y\in N^-(c)
    \quad\Longrightarrow\quad
        r_yU_z^nf-U_z^nf\in\mathrm{valIn}_A(k).
\end{equation}
The difference is measured in the ambient function module;
right translation by $y$ is not assumed to preserve $M_A$.
\begin{proof}
Left $\Gamma$-invariance, the $N$ and $T$ laws, and the
away action are preserved by OST5.  For continuity modulo
$\mathfrak m^k$, choose a depth at which both the character
and $f$ are trivial on the respective congruence subgroups.
The reduction of $f$ in $(A/\mathfrak m^k)^d$ is then
$I_1(c')$-fixed.  OST6 makes its image under $U_a$
$I_1(c')$-fixed as well.  Finite sums commute with reduction,
so the required uniform continuity lifts back to $A$.

For the second assertion, choose $c_1$ such that $f$ is
$N^-(c_1)$-invariant modulo $\mathfrak m^k$.  For $n\geq c_1$,
its reduction is fixed by $N,T_c(c),N^-(c+n)$.
Apply the last assertion of OST6 $n$ times.  The reduction of
$U_z^nf$ is fixed by $N^-(c)$, proving the congruence.
The threshold depends on $f,k,c$, as permitted by the source.
\end{proof}
% END OST7

% BEGIN OST8
\paragraph{OST8. Separated completeness and the initial value-filtration rules.}
\label{claim:OST8}
The modules $V_A(k)$ decrease with $k$ and satisfy
\[
 \mathfrak m^kM_A\subset V_A(k),\qquad
 \mathfrak m^kV_A(1)\subset V_A(k+1),\qquad
 \bigcap_kV_A(k)=0.
\]
If $f_{k+1}-f_k\in V_A(k)$, there is a unique $f\in M_A$
with $f-f_k\in V_A(k)$ for all $k$.
\begin{proof}
The inclusions follow by multiplying coordinates in the relevant
ideal powers.  Separation follows from separation of $A$.
For existence, take the $\mathfrak m$-adic limit of each
coordinate $f_k(g)_i$.  For every fixed precision the convergence
is uniform over all $g,i$, because the difference hypothesis is
a value-filtration statement.  All exact linear transformation
laws in \eqref{eq:ost-space} pass to the limit; the matrices and
scalar characters preserve the ideal powers.  For the continuity
law at precision $k$, use the depth supplied by $f_k$ and write
\[
 r_yf-f=r_y(f-f_k)+(r_yf_k-f_k)-(f-f_k).
\]
Every term is zero modulo $\mathfrak m^k$ uniformly in $g,i$.
Thus the limit lies in $M_A$, and uniqueness again follows
from separation.
\end{proof}
% END OST8

% BEGIN OST9
\paragraph{OST9. Finite-coordinate approximation of late operator iterates.}
\label{claim:OST9}
For every $k$ there are a finite set $D_k$, an $A$-linear
evaluation map $v_k:M_A\to A^{D_k}$, and a section
$s_k:A^{D_k}\to M_A$, such that
\begin{align}
 f\in V_A(j)&\ \Longrightarrow\ v_k(f)\in(\mathfrak m^j)^{D_k},
 \label{eq:ost-evaluation-filtration}\\
 \forall f\ \exists n_0\ \forall n\geq n_0\quad
 U_z^nf-s_kv_k(U_z^nf)&\in V_A(k).
 \label{eq:ost-finite-approx}
\end{align}
\begin{proof}
Choose $c$ using \eqref{eq:ost-contchar}, choose $x_i$ using
OST1, and set $D_k=\{1,\ldots,s\}\times\{1,\ldots,d\}$.
Let $v_k$ evaluate at those coordinates and let $s_k$ be OST3.
The first inclusion is immediate.  For large $n$, OST7 and
OST3 say that the two terms in \eqref{eq:ost-finite-approx}
are invariant modulo $\mathfrak m^k$ under $N^-(c)$.
Their difference vanishes at every $x_i$.  OST4 gives the
desired value congruence.  Only these finitely many evaluations
have finite reductions; finiteness of the whole quotient
$M_A/V_A(k)$ is neither asserted nor needed.
\end{proof}
% END OST9

\subsection{A complete-ring ordinary projector}

% BEGIN OST10
\paragraph{OST10. The original ordinary module is periodic at each precision.}
\label{claim:OST10}
Writing $T_0=U_z|_{M_A}$, the original definition is equivalent to
\begin{equation}\label{eq:ost-periodic}
 f\in O_A\quad\Longleftrightarrow\quad
 f\in M_A\ \hbox{and}
 \forall k\ \exists q\geq1\quad T_0^qf-f\in V_A(k).
\end{equation}
Every $U_a$ with $a\in P$ preserves $O_A$.
\begin{proof}
If $T_0^qf-f\in V_A(k)$, the telescoping identity
\[
 T_0^{qr}f-f=\sum_{i=0}^{r-1}T_0^{qi}(T_0^qf-f)
\]
shows the same assertion for every positive multiple of $q$.
Thus the kernels in \eqref{eq:ost-ord-definition} form a
directed family: periods $q,q'$ are both replaced by $qq'$.
The sum of a directed family of submodules is its union,
which proves \eqref{eq:ost-periodic}.  OST5 gives
$U_aT_0=T_0U_a$ on $M_A$, and OST7 preserves the space and
its filtration.  Apply $U_a$ to each periodic congruence.
\end{proof}
% END OST10

% BEGIN OST11
\paragraph{OST11. Every vector is eventually periodic modulo every precision.}
\label{claim:OST11}
For $f\in M_A$ and $k\geq0$, there are $a\geq0$ and $q\geq1$
such that $T_0^{a+q}f-T_0^af\in V_A(k)$.
\begin{proof}
The ring $A/\mathfrak m^k$ is finite.  Indeed, finite generators
of $\mathfrak m$ show that each
$\mathfrak m^j/\mathfrak m^{j+1}$ is finite dimensional over
the finite residue field, and induction proves the assertion.
Use OST9 and examine $v_k(T_0^nf)$ modulo $\mathfrak m^k$
for $n\geq n_0$.  Two of these tuples are equal, at indices
$a<a+q$.  Their section values differ by an element of
$\mathfrak m^kM_A\subset V_A(k)$, since $s_k$ is linear
on a finite free module.  Subtracting the two approximation
congruences proves the assertion.  Subsequent iterates repeat
with period $q$ because $T_0$ preserves $V_A(k)$.
\end{proof}
% END OST11

% BEGIN OST12
\paragraph{OST12. One factorial-limit projector on the complete space.}
\label{claim:OST12}
There is a unique $A$-linear endomorphism $e_A$ of $M_A$
satisfying
\begin{equation}\label{eq:ost-factorial-limit}
 \forall f\in M_A\ \forall k\ \exists J\ \forall j\geq J
                 \quad T_0^{j!}f-e_Af\in V_A(k).
\end{equation}
It preserves every $V_A(k)$, has image exactly $O_A$, and
is the identity on $O_A$.  In particular $e_A^2=e_A$.
\begin{proof}
Fix $f,k$ and the eventual period $q$ and starting index $a$
from OST11.  Once $j!\geq a$ and $q\mid j!$, the residue
of $T_0^{j!}f$ modulo $V_A(k)$ is independent of $j$:
all those exponents lie in the same residue class modulo $q$
in the periodic tail.  Therefore the factorial sequence is
Cauchy in the value filtration.  To use exactly OST8, choose
increasing integers $J_k$ so that all factorial terms of
index at least $J_k$ agree modulo $V_A(k)$.  The sequence
$T_0^{J_k!}f$ then satisfies its successive-difference
condition.  Its limit has \eqref{eq:ost-factorial-limit}
for the entire factorial sequence, not merely this subsequence.
It is unique by separation.

For $f+g$ or $bf$, the finite sum or scalar multiple of the
limits has the same defining congruences, using the larger
of the finitely many thresholds involved.  Uniqueness proves
linearity.  If $f\in V_A(k)$, every iterate is in that
submodule, so the limit is also in it.  The eventual period
$q$ for $f$ at precision $k$ gives
$T_0^q e_Af-e_Af\in V_A(k)$ by comparing with one sufficiently
late factorial term.  Thus $e_Af\in O_A$ by OST10.
If $f\in O_A$, a period $q$ already holds starting at zero.
For $q\mid j!$, $T_0^{j!}f-f\in V_A(k)$; hence
$e_Af-f\in V_A(k)$ for every $k$, proving $e_Af=f$.
The map is chosen once before $f,k$; no uniform threshold
over all vectors is asserted.
\end{proof}
% END OST12

% BEGIN OST13
\paragraph{OST13. Finite-coordinate approximation of ordinary vectors.}
\label{claim:OST13}
The maps of OST9 also satisfy, for all $f\in O_A$,
\begin{equation}\label{eq:ost-ordinary-approx}
                   f-e_As_kv_k(f)\in V_A(k).
\end{equation}
\begin{proof}
Choose a factorial index large enough for both OST9 and OST12.
Put $g=T_0^{j!}f$.  Then $g-f\in V_A(k)$, and
$g-s_kv_k(g)\in V_A(k)$.  By
\eqref{eq:ost-evaluation-filtration},
$s_kv_k(g-f)\in\mathfrak m^kM_A\subset V_A(k)$.
The identity
\[
 f-s_kv_k(f)=(g-s_kv_k(g))-(g-f)+s_kv_k(g-f)
\]
puts its left side in $V_A(k)$.  Apply the filtration-preserving
projector and use $e_Af=f$.
\end{proof}
% END OST13

% BEGIN OST14
\paragraph{OST14. One fixed finite generating map makes every next correction.}
\label{claim:OST14}
Fix the data at precision one, let $D=D_1$, and define
$\sigma=e_As_1:A^D\to O_A$.
For every $j\geq0$ and $f\in O_A\cap V_A(j)$ there is
$b\in(\mathfrak m^j)^D$ such that
\begin{equation}\label{eq:ost-one-step}
                         f-\sigma b\in V_A(j+1).
\end{equation}
\begin{proof}
Use OST13 at precision $j+1$.  Writing the standard vectors
of $A^{D_{j+1}}$ as $\delta_p$, put
$w_p=e_As_{j+1}\delta_p$ and
$a_p=(v_{j+1}f)_p\in\mathfrak m^j$.
Then $f-\sum_pa_pw_p\in V_A(j+1)$.
At precision one, each $w_p$ satisfies
$w_p-\sigma v_1(w_p)\in V_A(1)$.
Set $b=\sum_pa_pv_1(w_p)$.  Its coordinates lie in
$\mathfrak m^j$, and
\[
 f-\sigma b=
 (f-\sum_pa_pw_p)+\sum_pa_p(w_p-\sigma v_1(w_p)).
\]
The last sum belongs to $\mathfrak m^jV_A(1)\subset V_A(j+1)$
by OST8.  This proves the claim without assuming the desired
reverse filtration inclusion in advance.
\end{proof}
% END OST14

% BEGIN OST15
\paragraph{OST15. Successive approximation over the whole coefficient ring.}
\label{claim:OST15}
For every $j$ and $f\in O_A\cap V_A(j)$ there is
$b\in(\mathfrak m^j)^D$ with $\sigma b=f$.
\begin{proof}
Start with $r_0=f$.  Apply OST14 repeatedly to choose
$b_n\in(\mathfrak m^{j+n})^D$ and
$r_{n+1}=r_n-\sigma b_n\in V_A(j+n+1)$.
The series $b=\sum_{n\geq0}b_n$ converges coordinatewise
in the finite free module $A^D$.  Its sum belongs to
$(\mathfrak m^j)^D$: this follows either from the definition
of the inverse-limit topology or by reducing every partial
sum modulo $\mathfrak m^j$.  For every $N$,
\[
 f-\sigma b=r_N+
          \sigma\Bigl(\sum_{n<N}b_n-b\Bigr).
\]
The first term is in $V_A(j+N)\subset V_A(N)$.
The second is in $\mathfrak m^NO_A\subset V_A(N)$,
since its coefficient vector belongs to $(\mathfrak m^N)^D$.
Separation proves $f=\sigma b$.  No completeness or flatness
of $O_A$ was presumed; only completeness of $A$ and the
already proved value-filtration separation were used.
\end{proof}
% END OST15

% BEGIN OST16
\paragraph{OST16. Finite generation and exact ordinary filtration.}
\label{claim:OST16}
The module $O_A$ is finite over $A$ and, for every $j\geq0$,
\begin{equation}\label{eq:ost-exact-filtration}
                     O_A\cap\mathrm{valIn}_A(j)=\mathfrak m^jO_A.
\end{equation}
\begin{proof}
The case $j=0$ of OST15 gives surjectivity of the fixed map
$\sigma:A^D\to O_A$, hence finite generation.  For general
$j$, the same result expresses every vector in the left side
of \eqref{eq:ost-exact-filtration} as
$\sum_{p\in D}b_p\sigma\delta_p$ with $b_p\in\mathfrak m^j$.
The reverse inclusion follows by multiplying function values.
This proves both directions on the literal ordinary functions.
\end{proof}
% END OST16

% BEGIN OST17
\paragraph{OST17. Finite freeness without a reducedness hypothesis.}
\label{claim:OST17}
The module $O_A$ is finite free over $A$.
\begin{proof}
Let $v=v_1|_{O_A}$ and $\Psi=\sigma v$.  By OST13 and OST16,
$f-\Psi f\in\mathfrak mO_A$ for every $f$.  Nakayama's
lemma gives surjectivity of $\Psi$.  For completeness, the
form of Nakayama used here follows by choosing finite generators
of a module $Q=\mathfrak mQ$, expressing them through a matrix
with entries in $\mathfrak m$, and multiplying by the adjugate
of $1$ minus that matrix; its determinant is a unit.  Apply
this to $Q=\operatorname{coker}\Psi$.

A surjective endomorphism of a finite module over a Noetherian
ring is injective.  Indeed, the increasing sequence of kernels
of its powers stabilizes.  If $x\in\ker\Psi$, choose $y$
with $\Psi^ny=x$ at a stabilization index $n$; then
$y\in\ker\Psi^{n+1}=\ker\Psi^n$, giving $x=0$.
Thus $\Psi$ is invertible.  The maps
\[
 O_A\xrightarrow{v}A^D
       \xrightarrow{\Psi^{-1}\sigma}O_A
\]
have composite the identity.  Hence $O_A$ is a finite
projective module.  A finite projective module over a local
ring is free: lift a basis of its residue space to obtain
a surjection $A^r\to O_A$, split it, and call the kernel
$K$.  Reducing the direct sum modulo $\mathfrak m$ gives
$K/\mathfrak mK=0$, and finite Nakayama gives $K=0$.
The argument also covers $O_A=0$ and $d=0$.
\end{proof}
% END OST17

% BEGIN OST18
\paragraph{OST18. The distinguished ordinary operator is invertible.}
\label{claim:OST18}
The restriction of $U_z$ to $O_A$ is bijective.
\begin{proof}
It preserves $O_A$ by OST10.  For $f\in O_A$, choose its
period $q\geq1$ at precision one.  Then
$f-U_z(U_z^{q-1}f)\in V_A(1)\cap O_A=\mathfrak mO_A$.
Thus $U_z$ is surjective modulo $\mathfrak m$.  Apply finite
Nakayama to its cokernel, then the Noetherian kernel argument
of OST17.  This proves bijectivity over $A$, not only over
its residue field.
\end{proof}
% END OST18

% BEGIN OST19
\paragraph{OST19. Full original ordinary-structure endpoint.}
\label{claim:OST19}
For every original Frame $X$, every $A$ of class
\texttt{IsCNL}, every $d\geq0$, every $\tau,\theta,U$ with
\texttt{ContChar}, \texttt{FinLevel}, and \texttt{Free}, the
literal $X.\mathrm{Ord}(\tau,\theta,U)$ is finite free,
$X.U(z)$ is bijective on it, and its value filtration is its
intrinsic $\mathfrak m$-adic filtration.  There is one
$A$-linear map $e_A:X.\mathrm{Sp}\to X.\mathrm{Ord}$ which
fixes every ordinary function and satisfies
\eqref{eq:ost-factorial-limit} with the original
$X.U(z)$ and $X.\mathrm{valIn}$.
\begin{proof}
The spaces, maps, actions, and filtrations were identified
literally in \eqref{eq:ost-space}--\eqref{eq:ost-ord-definition}.
OST1--OST9 produce every input to the abstract construction
from the three stated hypotheses on that same Frame.
OST10--OST12 identify the projector image with its actual
ordinary submodule.  OST16--OST18 give the remaining
assertions.  Regarding $e_A$ as a map to its proved image
changes no function value.  These are exactly the fields
of \texttt{ordStruct}, including one projector before all
later vectors and precisions.
\end{proof}
% END OST19

% BEGIN OST20
\paragraph{OST20. Every positive-monoid operator is an ordinary unit.}
\label{claim:OST20}
For every $a\in P$, $U_a$ is bijective on $O_A$.  Together
with OST17 this is the full original \texttt{ordFree} endpoint.
\begin{proof}
Choose $h\geq0$ with $b=z^ha^{-1}\in P$ by the original
Iwahori axiom.  Since $a,b\in T_f$, they commute and
$ab=ba=z^h$.  OST5 gives, on $O_A$,
\[
                  U_aU_b=U_bU_a=U_z^h.
\]
The last map is invertible by OST18.  An explicit inverse of
$U_a$ is $U_bU_z^{-h}$: all these maps commute, and either
composite is the identity.  This includes $h=0$ and adds no
restriction on the residue characteristic.
\end{proof}
% END OST20

\subsection{The original coefficientwise base-change theorem}

% BEGIN OST21
\paragraph{OST21. Local coefficient maps preserve every original condition.}
\label{claim:OST21}
Let $\varphi:A\to B$ be any local homomorphism of complete
Noetherian local rings with finite residue fields.  Put
$\tau_B=\varphi(\tau)$ and $\theta_B=\varphi(\theta)$
entrywise.  The coefficientwise map
\[
 C_\varphi:\mathcal F_A\longrightarrow\mathcal F_B,
                 \qquad C_\varphi(f)(g)_i=\varphi(f(g)_i)
\]
restricts to $M_A\to M_B$ and $j:O_A\to O_B$.
The same finite-level and stabilizer hypotheses hold for $B$,
and $\theta_B$ is continuous in the required sense.
\begin{proof}
Locality gives $\varphi(\mathfrak m_A^k)\subset\mathfrak m_B^k$
for every $k$.  The exact transformation laws pass through
the ring map, and every value congruence maps to its counterpart
over $B$.  Finite sums in OST5 give
$C_\varphi U_a=U_a C_\varphi$.  Thus the periodic description
OST10 maps into the corresponding description over $B$.
The finite-level hypothesis contains no coefficient ring.
In the stabilizer predicate, the geometric assertions
$b\in Z(G_p)\cap T$ are unchanged, and applying $\varphi$
entrywise to $\theta(b)\tau(u)=1$ gives its exact mapped
version.  Applying $\varphi$ to
\eqref{eq:ost-contchar} proves continuity of $\theta_B$.
This proof does not assume that $\varphi$ is injective,
surjective, flat, or finite.
\end{proof}
% END OST21

% BEGIN OST22
\paragraph{OST22. One finite family gives simultaneous residual reconstruction.}
\label{claim:OST22}
Choose a depth $c$ at which $\theta$ is trivial modulo
$\mathfrak m_A$, and use the representatives $x_i$ and
the section $s_c$ of OST1 and OST3.  For
$p=(i,j)$ let $s_p=s_c\delta_p$ and $w_p=e_A s_p\in O_A$.
For $f\in O_A$ and $g\in O_B$ one has respectively
\begin{align}
 f-\sum_p f(x_i)_j w_p&\in\mathfrak m_AO_A,\label{eq:ost-bc-A}\\
 g-\sum_p g(x_i)_j j(w_p)&\in\mathfrak m_BO_B.
 \label{eq:ost-bc-B}
\end{align}
The evaluations in the two rings satisfy
$\operatorname{ev}_{B,p}(j(f))=\varphi(\operatorname{ev}_{A,p}(f))$.
\begin{proof}
Equation \eqref{eq:ost-bc-A} is OST13 at precision one
followed by OST16.  The evaluation compatibility is the
definition of $j$.  To prove \eqref{eq:ost-bc-B}, set
\[
 h=\sum_p g(x_i)_j C_\varphi(s_p),\qquad
 g'=\sum_p g(x_i)_j j(w_p)\in O_B.
\]
All $C_\varphi(s_p)$ lie in $M_B$ and are
$N^-(c)$-invariant modulo $\mathfrak m_B$ by OST3 and
OST21.  They have the same delta-function evaluations.
By OST12 over $B$, $U_z^{n!}g\equiv g$ modulo the
value filtration for sufficiently large $n$.  By OST7
the same late iterate is $N^-(c)$-invariant modulo
$\mathfrak m_B$.  Hence $g$ itself has that invariance.
Since $h$ agrees with $g$ at each $x_i$, OST4 gives
$g-h\in\mathrm{valIn}_B(1)$.

For each of the finitely many $p$, OST12 over $A$ gives
$U_z^{n!}s_p-w_p\in\mathrm{valIn}_A(1)$ for all
sufficiently large $n$.  Choose one $n$ beyond those
finitely many thresholds and the threshold for $g$.
Naturality of the finite sums and OST21 imply
$U_z^{n!}h-g'\in\mathrm{valIn}_B(1)$.
Applying $U_z^{n!}$ to $g-h$, and then subtracting the
two congruences, proves $g-g'\in\mathrm{valIn}_B(1)$.
Both terms are ordinary, so OST16 over $B$ converts this
into \eqref{eq:ost-bc-B}.  This argument uses the same
$A$-valued sections and ordinary vectors in both rings;
an unrelated basis of $O_B$ would not prove the claim.
\end{proof}
% END OST22

% BEGIN OST23
\paragraph{OST23. A finite-module base-change criterion with an explicit inverse.}
\label{claim:OST23}
Suppose $O$ and $O'$ are finite modules over local Noetherian
rings $A$ and $B$, $j:O\to O'$ is $A$-linear, and finite
families $w_p\in O$, $A$-linear maps $v_p:O\to A$, and
$B$-linear maps $v'_p:O'\to B$ satisfy the evaluation
compatibility and the two congruences
of OST22.  Then
\[
 L:B\otimes_A O\longrightarrow O',\qquad b\otimes f\longmapsto b\,j(f)
\]
is an isomorphism.
\begin{proof}
Let $\Psi(f)=\sum_p v_p(f)w_p$.  Its reduction modulo
the maximal ideal is the identity.  The finite Nakayama and
Noetherian kernel arguments in OST17 make $\Psi$ an
automorphism.  Define
\[
 R:O'\longrightarrow B\otimes_AO,\qquad
         R(g)=\sum_p v'_p(g)(1\otimes w_p).
\]
On a pure tensor the compatibility of evaluations gives
$RL=1\otimes\Psi$, an automorphism, and thus $L$ is
injective.  The second congruence says
$g-LR(g)\in\mathfrak m_BO'$.  The finite module
$\operatorname{coker}L$ is therefore equal to its maximal-ideal
multiple and vanishes by Nakayama.  Hence $L$ is surjective.
Its inverse is $(1\otimes\Psi)^{-1}R$.
No flatness of $B$ over $A$ was used; injectivity comes from
this actual left inverse.
\end{proof}
% END OST23

% BEGIN OST24
\paragraph{OST24. Full original coefficientwise ordinary base change.}
\label{claim:OST24}
Under the original hypotheses of OST19 and every local map
$A\to B$ between rings of class \texttt{IsCNL}, the
coefficientwise map of OST21 induces an isomorphism
\begin{equation}\label{eq:ost-full-base-change}
 B\otimes_A X.\mathrm{Ord}_A(\tau,\theta,U)
       \ \xrightarrow{\ \sim\ }
 X.\mathrm{Ord}_B(\varphi\tau,\varphi\theta,U),
 \qquad b\otimes f\longmapsto\bigl(g\mapsto b\varphi(f(g))\bigr).
\end{equation}
This is the exact \texttt{IsBaseChange B j} assertion of
\texttt{OrdBCMapOrd.ordBC}, with its stipulated formula for $j$.
\begin{proof}
OST21 transports every input hypothesis to the same Frame
and level over $B$.  OST19 gives finite modules on both sides.
OST22 supplies the evaluation-compatible data of OST23, which
proves \eqref{eq:ost-full-base-change}.  Its map is the
tensor extension of the original coefficientwise map, by the
displayed pure-tensor formula.  This identifies the map, not
merely its source and target ranks.  The proof includes
nonreduced rings and maps with kernels.
\end{proof}
% END OST24

% BEGIN OST25
\paragraph{OST25. The projector commutes with coefficient change and positive operators.}
\label{claim:OST25}
For every coefficient map in OST21 and every $a\in P$,
\[
 C_\varphi e_A=e_B C_\varphi\quad\hbox{on }M_A,
                  \qquad U_ae_A=e_AU_a.
\]
\begin{proof}
The first equality follows by applying $C_\varphi$ to
\eqref{eq:ost-factorial-limit}.  It preserves precision $k$
by locality and commutes with $U_z^{j!}$ by OST5, so both
sides are the unique factorial limit of the same function
over $B$.  For the second equality, $U_a$ preserves the
value filtration and commutes with $U_z$ on $M_A$ by OST5
and OST7.  Apply it to the same limit and use uniqueness.
These equalities also show that the base-change isomorphism
of OST24 intertwines the original positive-monoid operators.
\end{proof}
% END OST25

\subsection{Installing the actual local matrix datum}

% BEGIN OST26
\paragraph{OST26. Every axiom of the original local general-linear Iwahori datum.}
\label{claim:OST26}
Let $K$ be any of the original nonarchimedean local fields, with
valuation ring $\mathcal O$, maximal ideal $\mathfrak p$, and the
source's chosen uniformizer $\pi=\texttt{unif}(K)$.  Normalize the
additive valuation by $v(\pi)=1$.  For rank $n\geq0$, the literal
\texttt{glData} is
\begin{align*}
 N_f&=\{\text{upper unipotent matrices over }K\},\qquad
 T_f=\{\text{invertible diagonal matrices over }K\},\\
 N&=N_f\cap\operatorname{GL}_n(\mathcal O),\qquad
 T=\{\operatorname{diag}(u_i):u_i\in\mathcal O^\times\},\\
 T_c(c)&=\{\operatorname{diag}(u_i):u_i\equiv1
                                  \pmod{\mathfrak p^{c+1}}\},\\
 N^-(c)&=\{\text{lower unipotent matrices congruent to }1
                                  \pmod{\mathfrak p^{c+1}}\},\\
 P&=\{\operatorname{diag}(\pi^{a_i}):a_0\geq\cdots\geq a_{n-1},
                                      \ a_i\in\mathbb Z\},\\
 z&=\operatorname{diag}(\pi^{n-1-i})_{0\leq i<n}.
\end{align*}
These actual groups satisfy all the original \texttt{IwaData.Ax}
fields, so they supply the exact \texttt{glIwa} used above.
\begin{proof}
The source's $\mathrm{Kc}(r)$ is the kernel of reduction from
$\operatorname{GL}_n(\mathcal O)$ modulo $\mathfrak p^r$, mapped
into $\operatorname{GL}_n(K)$.  The integral map is injective, so
this is precisely the displayed congruence condition.  Inverse
transpose carries upper unipotents to lower unipotents, proving
the stated interpretation of $N^-(c)$.  Diagonal units identify
the source's valuation-one torus with $T$.  The selected
$\texttt{unif}(K)$ has uniformizer valuation by its defining
choice; no different uniformizer or monoid is substituted.

The inclusions, decreasing congruence subgroups, commutativity
of $T_f$, and membership of $z$ in $P$ follow directly.
Conjugation by $\operatorname{diag}(t_i)$ multiplies the
$(i,j)$ entry by $t_i/t_j$.  It therefore preserves upper
unipotents, and when the $t_i$ are units it preserves the
integral upper unipotents and every lower congruence group.
For $a=\operatorname{diag}(\pi^{a_i})\in P$, the difference
$a_i-a_j$ is nonnegative when $i<j$, proving
$aNa^{-1}\subset N$.  For $i>j$, conjugation by $a^{-1}$
multiplies the lower entry by $\pi^{a_j-a_i}$ with
nonnegative exponent, proving
$a^{-1}N^-(c)a\subset N^-(c)$.  For $a=z$ this exponent
equals $i-j\geq1$, which gives the required strict inclusion
into $N^-(c+1)$.

For the swap axiom, $y\in N^-(c)$ and $x\in N$ imply
$yx\equiv x\pmod{\mathfrak p^{c+1}}$.  We give the integral
upper-diagonal-lower elimination that applies to any matrix
$g$ with upper-unipotent reduction modulo $\mathfrak p^r$,
$r\geq1$.  Write
\[
 g=\begin{pmatrix}A&b\\c_0&d_0\end{pmatrix}
   =\begin{pmatrix}1&bd_0^{-1}\\0&1\end{pmatrix}
    \begin{pmatrix}A-bd_0^{-1}c_0&0\\0&d_0\end{pmatrix}
    \begin{pmatrix}1&0\\d_0^{-1}c_0&1\end{pmatrix}.
\]
Here $d_0\equiv1$ is a unit, $c_0\equiv0$, and the Schur
complement again has upper-unipotent reduction.  Induction on
$n$ yields $g=x'ty'$ with $x'\in N$,
$t\in T_c(r-1)$, and $y'\in N^-(r-1)$.  Apply this with
$r=c+1$.  For uniqueness, if an upper-unipotent matrix times
a diagonal matrix times a lower-unipotent matrix is $1$,
then the first two factors give a matrix that is simultaneously
upper triangular and lower unipotent.  It must be diagonal
with diagonal $1$, so all three factors are $1$.  This proves
the exact stronger $N_f,T_f,N^-(0)$ uniqueness field.

All required indices are finite.  For $a\in P$, choose
$D\geq0$ at least every $a_i-a_j$ with $i<j$.
The kernel of reduction of $N$ modulo $\mathfrak p^D$
is contained in $aNa^{-1}$, by the entrywise conjugation
formula.  Since $\mathcal O/\mathfrak p^D$ is finite, this
gives $[N:aNa^{-1}]<\infty$.  The same reduction argument
gives $[N^-(0):N^-(c)]<\infty$, while reduction of diagonal
units gives $[T:T_c(c)]<\infty$.  When $D=0$ the quotient
has one element; no positive lower bound for $D$ is needed.

For the exhaustion field, conjugating an upper-unipotent
matrix $x$ by $z^h$ multiplies its $(i,j)$ entry, $i<j$,
by $\pi^{h(j-i)}$.  A large enough $h$ makes the finitely
many entries integral, and an integral unipotent matrix has
integral inverse by the finite nilpotent geometric series.
Thus $z^h xz^{-h}\in N$.  For the positive-monoid domination
field, choose
\[
          h\geq\max_{0\leq i<n-1}(a_i-a_{i+1}),\qquad h\geq0.
\]
Then $h(n-1-i)-a_i$ is decreasing, so $z^ha^{-1}\in P$.
For the final torus-generation field, write each diagonal
entry $t_i=u_i\pi^{d_i}$ with $u_i\in\mathcal O^\times$.
Choose $h\geq0$ so that $d_i+h(n-1-i)$ is decreasing.
Then
\[
 \operatorname{diag}(t_i)
 =\operatorname{diag}(u_i)\,
   \operatorname{diag}(\pi^{d_i+h(n-1-i)})\,(z^h)^{-1},
\]
which is exactly $s a b^{-1}$ with $s\in T$, $a,b\in P$.
For ranks zero and one all missing adjacent-index conditions
are empty; choose $h=0$.  These calculations account for every
field of the original record.
\end{proof}
% END OST26

% BEGIN OST27
\paragraph{OST27. The product datum and the original Frame are the same objects.}
\label{claim:OST27}
The finite product of original Iwahori data, with all subgroups
and $P$ taken coordinatewise and $z=(z_v)_v$, satisfies all
the original axioms.  In particular OST26 constructs the
literal $\mathrm{iwS}(n,S)$, hence the Iwahori field of the
original $H.X$ for $H=\mathrm{hset}(Q,\tau_0)$.
\begin{proof}
Inclusions, conjugation laws, decrease with $c$, uniqueness,
and the swap identity are coordinatewise.  Choose the factors
in each coordinate for the swap and assemble their product
vectors.  Conjugation of the product $N$ is the product of
the conjugated coordinate subgroups; a finite product of their
finite coset sets proves each finite-index condition.

The two fields with a natural exponent need a common exponent.
For exhaustion choose $h_v$ in every coordinate and take their
maximum $h$.  If an exponent works, every larger one works:
conjugate further by a power of $z_v\in P_v$, which preserves
$N_v$.  For domination use the same maximum, since
$z_v^h a_v^{-1}=z_v^{h-h_v}(z_v^{h_v}a_v^{-1})\in P_v$.
The torus-generation factors are chosen coordinatewise and
then assembled.  The empty product also satisfies these
statements, with exponent zero.

The definitions in \texttt{TcAnti}, \texttt{Distinct}, and
\texttt{CoefficientRepTrivialOnUnramified} take exactly these
coordinatewise groups to define $\mathrm{piIwa}$ and
$\mathrm{iwS}$.  In \texttt{IsTheta.hset} the index set is
the original finite set $V.\mathrm{frame}.V_\ell$ and the
local field in coordinate $v$ is the actual completion
$L_{s.w(v)}$.  Thus every group, congruence depth, operator
representative quotient, and selected uniformizer used in
OST1--OST25 is that of the original Frame, with no replacement
of the group by an abstractly isomorphic one.
\end{proof}
% END OST27

% BEGIN OST28
\paragraph{OST28. Full ordinary structure at every original arithmetic level.}
\label{claim:OST28}
Fix any original $\ell,L,\Omega,s,n\geq1,V,Q,c,d,\tau_0,r,\Lambda$
in the scope of OTW12, set $H=\mathrm{hset}(Q,\tau_0)$,
$A=\Lambda.\Lambda$, and fix the single character $\theta$
provided there before choosing a level.  For every original
compact-open auxiliary level $K$, the actual module
\[
 H.X.\mathrm{Ord}\bigl(H.\tau_A(A),\theta,H.\mathrm{lev}(K)\bigr)
\]
is finite free; its original positive-monoid operators are
bijective; it has the exact value filtration and one
factorial-limit projector of OST19.  For every local map
$A\to B$ with $B$ of class \texttt{IsCNL}, its ordinary
base change is exactly the original coefficientwise map.
\begin{proof}
OTW1 verifies \texttt{IsCNL} on this exact complete ring.
OST26--OST27 supply the actual Iwahori datum, including all
local matrix and finite-product axioms.  FNL10 supplies
\texttt{FinLevel} at each $H.\mathrm{lev}(K)$.  OTW12 gives
\texttt{ContChar} for its already chosen $\theta$, and its
full \texttt{IsWt} field gives \texttt{Free} for
$H.\tau_A(A)$ at every such level, using NTL17.
Apply OST19, OST20, and OST24 on these very objects.
The projector may depend on $K$, as the source allows, but
within each level it is chosen before all functions and
precisions.  The character is still fixed before every $K$.
Neither \texttt{TypeForAlt} nor any of the other
$H.\mathrm{Ax}$ fields is needed for these applications.
The proof does not introduce any coefficient threshold or
restriction on $\ell$, and retains the original rank-four
case $\ell=3$.
\end{proof}
% END OST28

\subsection{Source correspondence and the remaining arithmetic entrance}

The definitions above are those of
\texttt{ContinuousTorusCharacter.lean}: the Iwahori record at
lines 29--67, $\mathrm{ContChar}$ at 136--137, the Frame and
its finite-level and stabilizer predicates at 151--175, and
$\mathrm{Sp}$, $\mathrm{valIn}$, $\mathrm{Ord}$, and the
coefficient map at 199--237.  The precise line intervals are
recorded with file hashes in this chapter's source manifest.
The source endpoints are
\texttt{OrdFree.lean}, lines 295--437 and 441--481, and
\texttt{OrdBCMapOrd.lean}, lines 198--232.  The auxiliary
source statements are expanded into OST1--OST18 rather than
invoked as unproved mathematical inputs.

This completes the ordinary-structure and ordinary-base-change
part of roadmap M1 on the original arithmetic Frames in OST28.
The finite-level theorem supplies \texttt{FinLevel}; NTL17
and OTW12 supply the stabilizer and complete-character inputs.
The chapter introduces no new coefficient threshold, no choice
of a replacement ideal, and no restriction excluding residue
characteristic three.  Its construction is independent of the
missing integral-type theorem.  The ordinary maps are the
Frame's actual double-coset maps, distinct from the separate
local theta-Hecke construction.

The following remain downstream arithmetic tasks: the
integral local type, classical weight/control and automorphic
fibers, arithmetic-point separation, the four Hida endpoints,
and all-level automorphy on the same ideal.  None follows merely
from finite freeness or base change.  Thus no whole
\texttt{COUNT\_II\_SUPPORT} obligation is retired by this
chapter alone.
